% !TITLE 闻官军收河南河北
% !AUTHOR 杜甫
% !DYNASTY 唐
% !ORDER 27

\begin{poem}
剑外忽传收蓟北\textsuperscript{1}，初闻涕泪满衣裳。
却看妻子愁何在，漫卷诗书喜欲狂\textsuperscript{2}。
白日放歌须纵酒，青春作伴好还乡\textsuperscript{3}。
即从巴峡穿巫峡\textsuperscript{4}，便下襄阳向洛阳\textsuperscript{5}。
\end{poem}

\begin{annotations}
\notetext{1}{剑外忽传收蓟北：剑门关外忽然传来收复蓟北的消息。剑外，剑门关以南（今四川）；蓟北（jì），今河北北部，安史之乱叛军的根据地。}
\notetext{2}{漫卷诗书喜欲狂：胡乱地卷起诗书，高兴得快要发狂了。平日珍贵的诗书此刻也顾不上了。}
\notetext{3}{青春作伴好还乡：趁着春天的好时光，回到家乡去。青春，明媚的春光。}
\notetext{4}{巴峡穿巫峡：从巴峡穿过巫峡，沿着长江三峡顺流而下。}
\notetext{5}{襄阳向洛阳：襄阳在今湖北，洛阳在今河南。四个地名串联成归乡的路线图。}
\end{annotations}

\begin{appreciation}
《闻官军收河南河北》被后人称为杜甫的"生平第一快诗"。杜甫传世一千四百多首诗，写愁写苦的多，这一首，从头快到尾。

剑外忽传收蓟北，初闻涕泪满衣裳——剑外是剑门关以南，杜甫当时逃难在四川。蓟北是叛军的老巢。消息是"忽传"的，来得突然；眼泪是"初闻"就落下的。注意，这不是喜极而泣，是悲极而泣：安史之乱打了八年，杜甫被困在叛军占领的长安，逃难，漂泊，一路写下了《春望》那样的诗。"国破山河在，城春草木深"——那时他只能一个人偷偷掉眼泪。八年的眼泪堵在心里，此刻才找到出口。

却看妻子愁何在，漫卷诗书喜欲狂——却看，是下意识地回头看。人高兴到极点，第一反应不是自己跳起来，是先看家人：妻子的愁眉是不是展开了。漫卷，是胡乱卷起。杜甫是读书人，书比命还宝贝，此刻也顾不上了，随手一卷就扔到一边。"喜欲狂"三个字，在他笔下千载难逢。

白日放歌须纵酒，青春作伴好还乡——须，是非喝不可。青春在这里是春天的意思，阳光正好，正好赶路。还乡——他的家在洛阳，八年没回去了。即从巴峡穿巫峡，便下襄阳向洛阳——四个地名一口气报完，水路接旱路，一条回家的路线，在地图上画得清清楚楚。这一联连用四个动词，即、穿、下、向，一步比一步急，恨不得今天就跨进家门。

这首诗写于广德元年（763年），杜甫五十二岁，人在梓州。官军收复河南河北，意味着打了八年的安史之乱终于要结束了。从《采薇》里那个"我心伤悲，莫知我哀"的戍边战士，到杜甫，战争让人流了三千年的眼泪。而杜甫用这首诗告诉你：眼泪，也有另一种流法。

人一辈子真正高兴到发狂的时刻，不多。你以后要是碰上这样的大好事，别端着，像杜甫一样，该笑就笑，该哭就哭。
\end{appreciation}
